% Answers to Chapter 1's exercises, lectures/L01/appendix/b_exercises.md. The book's own answers:
% the lecture publishes none.
\answerchapter{1}{C Fundamentals}
Every exercise in this chapter is a Code exercise; what follows answers only the questions three of
them ask.

\answer{c1:ex:3}
It depends on the compiler and its flags. GCC's \code{-Wswitch} is part of \code{-Wall}, so with
\code{-Wall} a \code{switch} on a \code{state_t} that leaves out one of the named constants gives a
warning along the lines of ``enumeration value \code{STATE_FAST} not handled in switch''; a plain
\code{gcc main.c} without it stays silent. Clang enables \code{-Wswitch} by default. The warning
only appears while the \code{switch} has no \code{default} label: a \code{default} counts as
handling every remaining value, so it silences the check (\code{-Wswitch-enum} warns even then).
This is the help from the compiler that \secref{c1:a:3} gives as a reason to prefer an \code{enum}
over a block of \code{#define}d constants, whose \code{switch} no compiler can check.

\answer{c1:ex:5}
Its forward declaration goes in \file{temp_reading.c}, once, right after the \code{#include}s, not
in \file{temp_reading.h}: the helper is \code{static}, an implementation detail rather than part of
the module's public API. Its definition goes at the bottom of the file, after the public functions,
which are defined in the order the header declares them. A reader then meets the public API first
and reaches the private detail only by scrolling on. Following \secref{c1:a:2}, the declaration
leaves a by-value parameter unqualified and the definition makes it \code{const}.

\answer{c1:ex:7}
Because C passes every argument by value. A \code{temp_reading_t*} parameter is a copy of the
caller's pointer: the function can change the struct it points to, which is all
\code{temp_reading_reset()} needs, but setting the parameter itself to \code{NULL} changes only the
local copy, and the caller's pointer still holds the old address. To change the caller's pointer,
the function needs that pointer's address, a \code{temp_reading_t**}, and assigns \code{NULL}
through it, the way \code{set_to_null()} does in \secref{c1:a:7}.
